\documentclass[10pt,a4paper]{amsart}
\usepackage{amsmath,amssymb,amsthm,graphicx}
\usepackage{verbatim}
\newenvironment{reference}{\comment}{\endcomment}
%\newenvironment{reference}{}{}
\newenvironment{slogan}{\comment}{\endcomment}
\newenvironment{history}{\comment}{\endcomment}

% For commutative diagrams we use Xy-pic
\usepackage[all]{xy}

% We use 2cell for 2-commutative diagrams.
\xyoption{2cell}
% Theorem environments.
%
\theoremstyle{plain}
\newtheorem{theorem}[subsection]{Theorem}
\newtheorem{proposition}[subsection]{Proposition}
\newtheorem{lemma}[subsection]{Lemma}

\theoremstyle{definition}
\newtheorem{definition}[subsection]{Definition}
\newtheorem{example}[subsection]{Example}
\newtheorem{exercise}[subsection]{Exercise}
\newtheorem{situation}[subsection]{Situation}

\theoremstyle{remark}
\newtheorem{remark}[subsection]{Remark}
\newtheorem{remarks}[subsection]{Remarks}

\numberwithin{equation}{subsection}

% Macros
%
\def\lim{\mathop{\mathrm{lim}}\nolimits}
\def\colim{\mathop{\mathrm{colim}}\nolimits}
\def\Spec{\mathop{\mathrm{Spec}}}
\def\Hom{\mathop{\mathrm{Hom}}\nolimits}
\def\Ext{\mathop{\mathrm{Ext}}\nolimits}
\def\SheafHom{\mathop{\mathcal{H}\!\mathit{om}}\nolimits}
\def\SheafExt{\mathop{\mathcal{E}\!\mathit{xt}}\nolimits}
\def\Sch{\mathit{Sch}}
\def\Mor{\mathop{\mathrm{Mor}}\nolimits}
\def\Ob{\mathop{\mathrm{Ob}}\nolimits}
\def\Sh{\mathop{\mathit{Sh}}\nolimits}
\def\NL{\mathop{N\!L}\nolimits}
\def\CH{\mathop{\mathrm{CH}}\nolimits}
\def\proetale{{pro\text{-}\acute{e}tale}}
\def\etale{{\acute{e}tale}}
\def\QCoh{\mathit{QCoh}}
\def\Ker{\mathop{\mathrm{Ker}}}
\def\Im{\mathop{\mathrm{Im}}}
\def\Coker{\mathop{\mathrm{Coker}}}
\def\Coim{\mathop{\mathrm{Coim}}}

% Boxtimes
%
\DeclareMathSymbol{\boxtimes}{\mathbin}{AMSa}{"02}

%
% Macros for moduli stacks/spaces
%
\def\QCohstack{\mathcal{QC}\!\mathit{oh}}
\def\Cohstack{\mathcal{C}\!\mathit{oh}}
\def\Spacesstack{\mathcal{S}\!\mathit{paces}}
\def\Quotfunctor{\mathrm{Quot}}
\def\Hilbfunctor{\mathrm{Hilb}}
\def\Curvesstack{\mathcal{C}\!\mathit{urves}}
\def\Polarizedstack{\mathcal{P}\!\mathit{olarized}}
\def\Complexesstack{\mathcal{C}\!\mathit{omplexes}}
% \Pic is the operator that assigns to X its picard group, usage \Pic(X)
% \Picardstack_{X/B} denotes the Picard stack of X over B
% \Picardfunctor_{X/B} denotes the Picard functor of X over B
\def\Pic{\mathop{\mathrm{Pic}}\nolimits}
\def\Picardstack{\mathcal{P}\!\mathit{ic}}
\def\Picardfunctor{\mathrm{Pic}}
\def\Deformationcategory{\mathcal{D}\!\mathit{ef}}
\usepackage{cleveref}
\crefname{originallemma}{original lemma}{original lemmas}
\Crefname{originallemma}{Original lemma}{Original lemmas}
\makeatletter
\newlabel{lemma-symbol-when-one-is-a-unit}{{02Q2}{source}{}{sourcepointer-1}{}}
\newlabel{lemma-symbol-when-one-is-a-unit@cref}{{[originallemma][0][]02Q2}{[1][0][]source}{}{}{}}
\makeatother
\crefname{originallemma}{original lemma}{original lemmas}
\Crefname{originallemma}{Original lemma}{Original lemmas}
\makeatletter
\newlabel{lemma-multiplicativity-symbol}{{02Q0}{source}{}{sourcepointer-2}{}}
\newlabel{lemma-multiplicativity-symbol@cref}{{[originallemma][0][]02Q0}{[1][0][]source}{}{}{}}
\makeatother
\crefname{originallemma}{original lemma}{original lemmas}
\Crefname{originallemma}{Original lemma}{Original lemmas}
\makeatletter
\newlabel{algebra-lemma-ord-additive}{{02MC}{source}{}{sourcepointer-3}{}}
\newlabel{algebra-lemma-ord-additive@cref}{{[originallemma][0][]02MC}{[1][0][]source}{}{}{}}
\makeatother
\crefname{originalproposition}{original proposition}{original propositions}
\Crefname{originalproposition}{Original proposition}{Original propositions}
\makeatletter
\newlabel{proposition-length-determinant-periodic-complex}{{02QD}{source}{}{sourcepointer-4}{}}
\newlabel{proposition-length-determinant-periodic-complex@cref}{{[originalproposition][0][]02QD}{[1][0][]source}{}{}{}}
\makeatother
\title{The sum of the orders of the tame symbols over the height one primes of a two dimensional Noetherian local domain vanishes}
\author{The Stacks Project Authors}
\date{stacks-project:chow.tex:local-snapshot-20261012; GNU FDL 1.2 or later}
\begin{document}
\maketitle
\noindent\textit{Editorial note.} All delivered mathematical text is the author's own, copied from the pinned Stacks Project chapter file chow.tex; no statement or proof step was rewritten, reordered or removed.\par\medskip
\noindent\textit{Editorial note.} Preamble provenance: the pinned chapter file opens with an input of the shared preamble and carries no preamble text of its own, so the project's single Stacks-wide preamble staged at sources/stacks-chow/preamble.tex is excerpted verbatim as two spans, lines 14-18 (Xy-pic setup) and lines 154-217 (the native theorem declarations and mathematical macros). Nothing was added, removed or redefined; the wrapper class and the amsmath/amssymb/amsthm/graphicx base packages are supplied by the assembler.\par\medskip
\noindent\textit{Editorial note.} The counted statement opens with a reference block recording that for an excellent ring the result is Kato's Milnor K-theory proposition; the block is retained verbatim inside the counted span, and the chapter preamble carries the reference environment as a comment environment so it does not print.\par\medskip
\noindent\textit{Editorial note.} Canonical location of the counted proposition: Stacks tag 02QJ (sources/stacks-chow/pinned-original-source.tex line 18340).\par\medskip
\noindent\textit{Allowed standard prerequisites (editorial).}\begin{itemize}
\item Additivity of the tame symbol and additivity of the order functions. Same chapter, Lemma labelled lemma-multiplicativity-symbol (source line 17434) and Algebra, Lemma labelled algebra-lemma-ord-additive (Stacks tag 02MC), invoked at proof lines 18370-18373. Supply the reduction of the general case f, g in K\textasciicircum{}* to the case f = a, g = b in A, which is the first step of the proof.
\item The determinant of a two-periodic complex and its equality with the order sum. Same chapter, Proposition labelled proposition-length-determinant-periodic-complex (source line 18024), invoked at proof line 18381. Supplies the equivalence of the displayed order sum with the vanishing of e\_A(A/(ab), a, b); it is retained only as a tagged navigation pointer because its interval lies inside a public item.
\item Definition of the order function on the residue field of a one dimensional quotient. Algebra, Definition labelled algebra-definition-ord (Stacks tag 02MD), used implicitly through the order functions appearing in the statement. Fixes the meaning of ord\_\{A/q\} in both displayed sums.
\end{itemize}\medskip
\ifdefined\hypertarget\hypertarget{sourcepointer-1}{}\fi
\noindent\textit{Source reference.} Original lemma 02Q2: sources/stacks-chow/pinned-original-source.tex, source lines 17483--17483. Same-chapter pointer into the author environment lemma at source line 17483; invoked at proof line None of the counted proof. It is retained only as a navigation pointer; every load-bearing step is inside the delivered span.\par\medskip
\ifdefined\hypertarget\hypertarget{sourcepointer-2}{}\fi
\noindent\textit{Source reference.} Original lemma 02Q0: sources/stacks-chow/pinned-original-source.tex, source lines 17434--17434. Same-chapter pointer into the author environment lemma at source line 17434; invoked at proof line 18371 of the counted proof. It is retained only as a navigation pointer; every load-bearing step is inside the delivered span.\par\medskip
\ifdefined\hypertarget\hypertarget{sourcepointer-3}{}\fi
\noindent\textit{Source reference.} Original lemma 02MC: sources/stacks-chow/tags.txt, source lines 3439--3439. Cross-chapter standard prerequisite located in the project tag list staged in the local snapshot; invoked at proof line 18373 of the counted proof. It is retained only as a navigation pointer; every load-bearing step is inside the delivered span.\par\medskip
\ifdefined\hypertarget\hypertarget{sourcepointer-4}{}\fi
\noindent\textit{Source reference.} Original proposition 02QD: sources/stacks-chow/pinned-original-source.tex, source lines 18024--18024. Same-chapter pointer into the author environment proposition at source line 18024; invoked at proof line 18381 of the counted proof. It is retained only as a navigation pointer; every load-bearing step is inside the delivered span.\par\medskip
\begin{lemma}[Key Lemma]
\label{lemma-secondary-ramification}
\begin{reference}
When $A$ is an excellent ring this is \cite[Proposition 1]{Kato-Milnor-K}.
\end{reference}
Let $A$ be a $2$-dimensional Noetherian local domain with fraction field $K$.
Let $f, g \in K^*$.
Let $\mathfrak q_1, \ldots, \mathfrak q_t$ be the height
$1$ primes $\mathfrak q$ of $A$ such that either $f$ or $g$ is not an
element of $A^*_{\mathfrak q}$.
Then we have
$$
\sum\nolimits_{i = 1, \ldots, t}
\text{ord}_{A/\mathfrak q_i}(d_{A_{\mathfrak q_i}}(f, g))
=
0
$$
We can also write this as
$$
\sum\nolimits_{\text{height}(\mathfrak q) = 1}
\text{ord}_{A/\mathfrak q}(d_{A_{\mathfrak q}}(f, g))
=
0
$$
since at any height one prime $\mathfrak q$
of $A$ where $f, g \in A^*_{\mathfrak q}$
we have $d_{A_{\mathfrak q}}(f, g) = 1$ by
Lemma \ref{lemma-symbol-when-one-is-a-unit}.
\end{lemma}

\begin{proof}
Since the tame symbols $d_{A_{\mathfrak q}}(f, g)$ are additive
(Lemma \ref{lemma-multiplicativity-symbol}) and the order
functions $\text{ord}_{A/\mathfrak q}$
are additive (Algebra, Lemma \ref{algebra-lemma-ord-additive})
it suffices to prove the formula when $f = a \in A$ and
$g = b \in A$. In this case we see that we have to show
$$
\sum\nolimits_{\text{height}(\mathfrak q) = 1}
\text{ord}_{A/\mathfrak q}(\det\nolimits_\kappa(A_{\mathfrak q}/(ab), a, b))
= 0
$$
By Proposition \ref{proposition-length-determinant-periodic-complex}
this is equivalent to showing that
$$
e_A(A/(ab), a, b) = 0.
$$
Since the complex
$A/(ab) \xrightarrow{a} A/(ab) \xrightarrow{b} A/(ab) \xrightarrow{a} A/(ab)$
is exact we win.
\end{proof}

\end{document}
